\clearpage
\phantomsection\addcontentsline{toc}{chapter}{来源对应与范围}
\markboth{来源对应与范围}{}
\begin{center}{\bfseries 来源对应与范围}\end{center}
所有来源均在MIT OCW 8.02 Spring 2007网页实际链接集合中. 本表以文件ID定位, 不把旧页脚课号与网页课次强行合一. 完整URL、SHA256、大小、PDF页数与许可见随仓库提供的来源清单. 每个摘要的实质主题按原PDF读取; 旧日期及错配详见独立来源审查. 以下ID省去共同前缀summary\_, 对应原文件summary\_wXXdY.pdf.
\begin{longtable}{p{.13\linewidth}p{.68\linewidth}p{.10\linewidth}}
\toprule 摘要ID & 实际物理主题 & 正文章\\\midrule\endhead
w01d1 & 标量场/向量场及三种图示；重力场、电场；逆平方律和场对试探物体的力 & 1\\
w01d2 & 电荷、库仑场、叠加；线/面/体电荷密度；电偶极矩、力矩及非均匀场受力的定性认识 & 1--3\\
w01d3 & 开/闭曲线和曲面；标量与向量线积分；面积分、通量；以 Maxwell 积分式提示用途。Topics 虽写坐标/梯度，正文主要为积分 & 2\\
w02d1 & 组内学习与可视化；势能、外力作功；势差线积分、势的参考点；U=qV、$\ee=-\nabla V$；实验可视化简介 & 3\\
w02d2 & Gauss 定律、通量、闭合面、球/柱/面三类对称性、高斯面解题步骤 & 4\\
w02d3 & 连续电荷场积分、叠加、$\lambda/\sigma/\rho$ 与 dq & 2--3\\
w03d2 & 导体屏蔽、电容串并联、介质降低场并增大电容、介质中的 Gauss 式；Faraday ice pail 简介 & 5\\
w03d3 & Gauss 定律练习方法；三种对称性与实体/壳层两种分布；通量与包围电荷 & 4\\
w04d1 & 静电平衡导体内部零场；电容定义和平行板推导；充电储能；电场能密度 & 5\\
w04d2 & 导体屏蔽；电流、电流密度、漂移速度、微观/宏观 Ohm 定律及电阻率；正文提及介质而未展开串并联计算 & 5--6\\
w04d3 & 电容计算三步；平行板电容；用场能密度、充电作功、$Q^2/(2C)$三法算储能 & 5\\
w05d1 & RC 电路：支路/节点、守恒、电容充放电、初始/稳态直观、指数和$\tau=RC$；实验计量简介 & 6\\
w05d2 & 考试一信息；离散/连续电场、场线、叠加、点电荷势/组装能、E与V、Gauss、电容/介质、概念题能力清单 & 1--6\\
w06d1 & 磁偶极/条形磁铁/地磁；无磁单极的定性描述；运动电荷的 Lorentz 力与弯曲轨迹 & 7--8\\
w06d2 & Ampère 定律、圆环与不包围电流的环路；高对称性解题法；条形磁铁/Helmholtz线圈实验简介 & 7\\
w06d3 & 载流线元的力、电流环磁矩和力矩；三种右手规则与力矩向量 & 8\\
w07d1 & 磁矩方向、力矩、非均匀磁场中的力、$U=-\vct m\cdot\bb$；电流环与条形磁铁类比 & 8\\
w07d2 & 磁悬浮稳定性、抗磁材料、磁矩在均匀/非均匀场中的力和力矩、能量判方向 & 8\\
w07d3 & Ampère 定律、圆柱壳/电流板/螺线管对称性、环路和穿过电流的分区计算步骤 & 7\\
w09d1 & Faraday 定律、磁通量、动生/时变感应电动势、改变B/A/角度三方式、Lenz定律 & 9\\
w09d2 & Faraday/Lenz、涡流制动、互感与变压器、匝数与输出电压关系 & 9\\
w09d3 & 下落线圈的感应电流与受力；Faraday/Lenz方法复述 & 9\\
w10d1 & 自感、反电动势；电感初始/稳态行为；磁储能/能密度；RL电路微分式与指数、$\tau=L/R$ & 9--10\\
w10d2 & 自感、LR电路、磁能和$\tau=L/R$ & 9--10\\
w11d2 & 受迫LRC电路、机械类比、R/L/C相位、谐振、相量与电抗 & 10\\
w11d3 & 受迫LRC：电阻/电抗/阻抗区别；谐振峰/相位；相量；平均功率和RMS & 10\\
w12d1 & LC自由振荡、无驱动LRC阻尼、机械/能量类比；后半还讲受迫振荡、相位、相量与谐振 & 10\\
w12d2 & Maxwell四式、位移电流；平面电磁波、横向关系、$E_0=cB_0$；Poynting矢量与电路能流 & 11\\
w12d3 & 位移电流与能流；充放电电容和电阻中的Poynting方向 & 11\\
w13d1 & 由Faraday与Ampère-Maxwell推导一维波动方程；光速、E/B关系；两个矩形环路的场图 & 11\\
w13d2 & Maxwell、平面波周期/波长/传播、场的横向约束；Poynting矢量及电路能流 & 11\\
w13d3 & Topics为EM Radiation；正文主要是位移电流与电容/电阻能流，并提示电磁波练习 & 11\\
w14d1 & 电磁辐射的产生；偏振、振荡带电面、四分之一波长天线、微波实验简介 & 11--12\\
w14d2 & 相长/相消、双缝干涉、单缝衍射、有限缝宽包络与条纹 & 12\\
w15d1 & 干涉、薄膜干涉的定性介绍、双缝、单缝衍射及干涉实验简介 & 12\\
w15d2 & 期末安排与12项能力清单（第1项另含a–d）；含Biot-Savart、电偶极、磁偶极、感应、DC/AC、波、干涉衍射与能流 & 1--23\\
\bottomrule\end{longtable}
第7、17章补入毕奥--萨伐尔、厚电流板等定量推导, 第11章补入电流面能量平衡、天线边界与正确传播图像. 这些是本稿展开内容, 不宣称每份短摘要原来就包含完整证明. 原摘要的实验背景在相关正文说明其物理目的; 实验操作、仪器软件和未公开题面不冒充本册已取得的材料.

习题按独立数学任务或独立物理情境计为34题组; 叶级小问按原Q/Step或显式(a)/(b)、(i)/(ii)计, 合并正文和答题页重复内容. 原编号同一问要求方向、数值和解释, 仍算一问且均须完成. 共有134个正式编号小问和35个Sample Exam练习小问, 共169问. 所有解答为AI独立解答, 未取得对应官方答案.
\begin{center}\begin{tabular}{lrrr}\toprule 原件 & 本册章 & 题组 & 小问\\\midrule
solving01.pdf & 13 & 7 & 17\\
solving02.pdf & 14 & 5 & 10\\
solving03.pdf & 15 & 2 & 15\\
solving04.pdf & 16 & 2 & 13\\
solving05.pdf & 17 & 3 & 22\\
solving06.pdf & 18 & 3 & 14\\
solving07.pdf & 19 & 1 & 10\\
solving08.pdf & 20 & 2 & 18\\
solving09.pdf & 21 & 6 & 20\\
solving10.pdf & 22 & 2 & 18\\
solving11.pdf & 23 & 1 & 12\\
\midrule 合计 & & 34 & 169\\\bottomrule\end{tabular}\end{center}
来源勘误公开处理: solving08原件缺Q5编号, 不增造缺问; solving05页脚沿用6, solving06沿用7, 后续类似, 文件ID保持稳定; solving10 Q11/Q12交叉回指错误、Q13放电导数正负冲突, 在题解中说明. 不把字体提取错位当成原公式, 涉及绝对值、角度、源极性和半径均核实际图页.

本轮未选14章长篇Course Notes、周一/周三小组题、实验指导、幻灯片和在线平台题. 每周书面作业被官方明确列为不向OCW用户提供. 本轮考试材料只有能力及日期说明和习题内练习, 不是已取得的正式测验/期末卷. 8.022仅核官方课程信息, 没有归档其题卷或把不同学期答案配给它. 这些边界不会用“全部OCW已完成”描述.
\begin{thebibliography}{9}
\bibitem{ocw} Massachusetts Institute of Technology. MIT OpenCourseWare, 8.02 Physics II: Electricity and Magnetism, Spring 2007. 课程材料: Peter Dourmashkin, Eric Hudson, Sen-Ben Liao及署名教师团队. \url{https://ocw.mit.edu/courses/8-02-physics-ii-electricity-and-magnetism-spring-2007/}.
\bibitem{read} MIT OCW. Readings: 36份本轮实际选中摘要. \url{https://ocw.mit.edu/courses/8-02-physics-ii-electricity-and-magnetism-spring-2007/pages/readings/}.
\bibitem{ps} MIT OCW. Problem Solving: 11份本轮实际选中周五习题课. \url{https://ocw.mit.edu/courses/8-02-physics-ii-electricity-and-magnetism-spring-2007/pages/problem-solving/}.
\bibitem{next} Gabriella Sciolla. 8.022 Physics II: Electricity and Magnetism, Fall 2004. 本轮仅核课程元数据, 后续阅读入口. \url{https://ocw.mit.edu/courses/8-022-physics-ii-electricity-and-magnetism-fall-2004/}.
\bibitem{license} MIT OCW. Terms and Conditions及第三方权利说明. \url{https://ocw.mit.edu/pages/privacy-and-terms-of-use/}; \url{https://mitocw.zendesk.com/hc/en-us/articles/4414756181403-How-is-all-rights-reserved-content-different-from-the-rest-of-OCW-content}.
\end{thebibliography}
